\subsection{Why does our system exist?}
In our research, we couldn't find a singular system that allows users to gain a general overview of public amenities, for example, parking, 
bike infrastructure, or public transport.
In Ireland and the United Kingdom, the prevalence of proper digitised records in county administrations varies wildly.
Some make use of state-of-the-art geographical information systems (GIS) while others can only rely on spreadsheets which are manually kept up to date.
\\ \\
A system that would allow users to quickly inspect a combined dataset grounded in automatically generated, real-world data could accelerate processes 
like planning permissions, urban development or resource allocation.
This applies especially to smaller counties that may not have the experienced personnel.
\\ \\
Improving how the government allocates resources could lead to improvements in efficiency, making public services more attainable, even for smaller communities.
For example, having accurate data on public transport and bike infrastructure could encourage more people to choose sustainable modes of transport, which would help reduce carbon and noise emissions. 
It could also make it easier for people in underserved communities to access important services, such as housing, helping to bridge social gaps.
\\ \\
Currently, solutions are fragmented, focusing on just one type of amenity, like transport or housing, without giving a full picture.
Even the more advanced GIS systems, while powerful, usually need a lot of manual input or additional programming, which smaller municipalities or organizations might not be able to handle.
Many of these GIS tools are expensive, or require specialised training, making them out of reach for smaller counties or individual users.

\subsection{The core technical problem}
The problem at the heart of our project is to find a way to source useful data about public amenities, without the need for any existing information or manual data entry.
This would allow the solution to be useful no matter the state of the users digital records.
Relying purely on existing data would do nothing to reduce the divide between counties with more data- and counties with less data.
This necessitates the creation of an interconnected series of modules that form a pipeline for the automatic extraction and fusion of data.
\\ \\
In order to present users with a map containing information about the number of public amenities, we must first identify them.
Our biggest technical challenge is the extraction of these data points. Different forms of public amenities exist in satellite
imagery, we can combine this with other data sources interweaving them, and producing a better result.
\\ \\
Detecting public amenities from overhead views can be tricky, as satellite imagery suffers from low resolution,
occlusion, seasonal changes, or cloud cover. This means classifying specific objects can be complicated, and can take a lot of time
and processing power. For instance, the task of quantifying on-street parking is dependent on reliably detecting cars that are not
actively participating in traffic (in aerial images). Not only that, but the detected points in the images must be mapped
back onto real-world coordinates to be of any use.
\\ \\
Ensuring consistency between our extracted data, and other data sets, can have its own pitfalls.
For example, compensating for:
\begin{itemize}
  \item{different precision in geographical alignment}
  \item{varied update frequencies}
  \item{different data formats}
\end{itemize}
Not only can data be geographically misaligned, but since satellite data is only a snapshot in time, it can lead to
temporal misalignment as well. Even the appearance of amenities can change depending on their location.
While cars look more or less the same anywhere on earth, the same cannot be said for public transport links, housing or hospitals,
which could lead to difficulties achieving the same level of accuracy in different regions.
Processing a large amount of data like this is no trivial task and finding a strategy to do this efficiently is necessary.

\subsection{A review of similar systems}
